% !TEX root = coursework_imu.tex
\documentclass[a4paper, 14pt]{extarticle}

\usepackage{fontspec}

% Основной шрифт — Times New Roman
\setmainfont{Times New Roman}
\setmonofont{Liberation Mono}


\usepackage{listings}
\lstset{
    basicstyle = \ttfamily,
    numbers = left,
    numbersep = 5pt,
    tabsize = 4,
    showstringspaces = false,
    keepspaces = true
}
\usepackage{amsmath}
\usepackage{enumitem}
\usepackage{indentfirst}
%\usepackage[T2A]{fontenc}
\usepackage[english, russian]{babel}
\usepackage{multirow}
\usepackage{tabularx}

\usepackage{subcaption}

\usepackage{array}
\usepackage{booktabs}
\usepackage{colortbl}
% Центрируем содержимое multirow по вертикали
\usepackage[table]{xcolor}
\usepackage{tikz}
\usetikzlibrary{backgrounds}

% удобный тип колонки с фиксированной шириной и центрированием
\newcolumntype{P}[1]{>{\centering\arraybackslash}p{#1}}

% малый корректор (если нужно подровнять край — поставьте 0.4pt или 0.2pt)
\newlength{\diagfudge}
\setlength{\diagfudge}{0pt}

% #1 = цвет нижне-левого треугольника
% #2 = цвет верхне-правого треугольника
% #3 = содержимое ячейки (одно значение)
\newcommand{\diagcell}[3]{%
  \begingroup
  \setlength{\fboxsep}{0pt}% убрать лишние отступы вокруг коробок
  % рисуем фон в "нулевой" коробке, сдвинутой влево на \tabcolsep
  \makebox[0pt][l]{\hspace*{-\tabcolsep}%
    \tikz[baseline=(box.base)]{%
      \node[inner sep=0pt, outer sep=0pt, anchor=west,
            minimum width=\dimexpr\linewidth+2\tabcolsep+\diagfudge\relax,
            minimum height=1.5\baselineskip] (box) {};
      \begin{scope}[on background layer]
        \fill[#1]  (box.south west) -- (box.south east) -- (box.north west) -- cycle;
        \fill[#2]  (box.north east) -- (box.south east) -- (box.north west) -- cycle;
        % Для отладки раскомментируйте следующую строку, чтобы увидеть рамку узла:
        % \draw[red,thin] (box.north west) rectangle (box.south east);
      \end{scope}
    }%
  }%
  % Теперь выводим само содержимое, ровно в ширине \linewidth и по центру
  \makebox[\linewidth][c]{\strut #3}%
  \endgroup
}


\usepackage{textcomp}
\newcommand{\degree}{\ensuremath{{}^\circ}} %использвоание значка градуса в $$

\usepackage{geometry}
\geometry{top=20mm}
\geometry{bottom=20mm}
\geometry{left=30mm}
\geometry{right=15mm}


\usepackage{titlesec}
\titleformat{\section}{\normalsize\it\centering}{\thesection}{1em}{}
\titleformat{\subsection}{\normalsize\it}{\thesubsection}{1em}{}

%\titleformat{\subsection}
%  {\normalfont\normalsize\bfseries\centering}
%  {\thesubsection}{1em}{}


\usepackage{tocloft}

%\renewcommand{\cftsecaftersnum}{.} % Добавляет точки после номера

% Заменяем лидеры на точки для всех уровней
\renewcommand{\cftdotsep}{1}
\renewcommand{\cftsecleader}{\cftdotfill{\cftdotsep}}
\renewcommand{\cftsubsecleader}{\cftdotfill{\cftdotsep}}
\renewcommand{\cftsubsubsecleader}{\cftdotfill{\cftdotsep}}

% Ваши текущие настройки
\tolerance=1000
\emergencystretch=3em
\pretolerance=150
\usepackage[none]{hyphenat}

% Добавьте эти строки для подавления предупреждений
\hbadness=10000  % не показывать предупреждения об underfull
\vbadness=10000  % то же для вертикального направления


\usepackage{graphicx}
\graphicspath{ {./figures/} }

\usepackage{caption}
\captionsetup[figure]{labelsep=endash}

\usepackage[
    %hidelinks  % ссылки без цвета и рамок, но остаются кликабельными
]{hyperref}


\usepackage{newtxmath}

\usepackage{caption}
\captionsetup[table]{justification=raggedright,singlelinecheck=false,labelsep=endash}

\usepackage{float}

